\chapter{Національна програма інформатизації}

Мета Національної програми інформатизації (НПІ) --- створення цифрового простору як
проміжного етапу розвитку інформаційного суспільства. Це прозоре правове середовище,
яке є надійно захищеним, базується на міжнародних стандартах, повноцінно забезпечує інформаційні потреби
та реалізує права і свободи громадян на основі своєчасної, достовірної та вичерпної інформації,
водночас істотно підвищуючи ефективність державного управління.

Національна програма інформатизації веде свою історію від Закону України № 74/98-ВР 1998 року,
який містив 28 статей, і до поточного Закону 2024 року № 2807-IX, що складається з 15 статей.
Головним чином НПІ закріплює протоколи запуску та термінації (завершення) програм, які охоплюють такі функції:
експертизу, аналіз, формування, контроль, виконання та звітування щодо програм інформатизації
як на державному (галузеві програми), так і на місцевому (самоврядування) рівнях. Вона також чітко визначає
суб'єктів керування: генерального замовника --- Міністерство цифрової трансформації,
Керівника (відповідальну посадову особу), а також виконавців та спеціалізованих підрядників.
Зокрема, НПІ є власником і розробником централізованої системи обліку таких програм. Крім того, Програма регламентує
процеси розробки апаратного та програмного забезпечення згідно з міжнародними стандартами ISO/IEC 12207 та ISO 9001,
а процеси супроводу та підтримки --- згідно з ISO/IEC/IEEE 14764:2022.

Основне завдання Національної програми інформатизації — безперервна оптимізація підприємств
у структурі органів виконавчої влади (ОВВ). Для досягнення цієї мети система державного управління повинна
мати інструменти для визначення протоколів інституційної самотрансформації, а також забезпечувати
узгоджену архітектуру програмних комплексів серед низки критично важливих міністерств.

\section{Загальні принципи}

На мета-рівні безперервний процес реформування органів виконавчої влади може бути описаний
як такий, що керується трьома фундаментальними правилами. У структурі кожного міністерства мають бути наявні:
\begin{enumerate}
    \item \textbf{Ін'єктивність управління юстиції} (як необхідний атрибут задля забезпечення внутрішнього аудиту та службових розслідувань);
    \item \textbf{ІТ-департамент або ІТ-управління} (яке виступає як окрема юридична особа з власним кодом ЄДРПОУ, стає продуктовим оунером (власником) для цифрових рішень Міністерства);
    \item \textbf{Управління трансформації} (підрозділ, що автоматизує та механізує організаційні процеси за допомогою ІТ-продуктів, розроблених ІТ-департаментом).
\end{enumerate}

Далі вже інші профільні управління можуть додаватися залежно від
специфічних функцій відповідного міністерства, але ці три стовпи, а також патронатна служба — є обов'язковими.
Міністерство юстиції, судова система, Генеральна прокуратура, поліція
та інші антикорупційні чи правоохоронні агенції, такі як НАЗК і НАБУ, повинні мати безпосередній або опосередкований
доступ до першого компонента (управління юстиції). Цей доступ має суворо регулюватися відповідними правилами ABAC (Attribute-Based Access Control).
Зі свого боку, Мінцифра, як координатор усього метапроцесу цифрової трансформації, отримує авторизований
доступ до третього компонента (управління трансформації). Також Мінцифра координує роботу та безпосередньо взаємодіє з
ІТ-департаментами кожного міністерства для забезпечення захищених каналів зв'язку до державних реєстрів, які
підтримуються силами ІТ-управлінь конкретних міністерств. Архітектура та шини обміну даними між усіма системами електронного документообігу
й реєстрами координуються ІТ-управлінням Мінцифри та платформою «Дія».

Завдання державної цифровізації та трансформації передбачає виконання таких стратегічних цілей:
\begin{itemize}
    \item кожне міністерство матиме свій автономний і висококомпетентний ІТ-департамент, який операційно обслуговуватиме відповідні реєстри;
    \item міністерства матимуть власні управління трансформації, що працюватимуть на базі єдиного продукту, у якому вони зможуть зручно моделювати свої бізнес- та адміністративні процеси;
    \item Мінцифра, як головний координатор, розроблятиме й затверджуватиме єдині регламенти роботи для ІТ-управлінь та управлінь трансформації кожного міністерства.
\end{itemize}
Платформа «Дія» при цьому матиме не лише базову шину документообігу та коробочне рішення «Дія: Документообіг»,
але й виступатиме провайдером ліцензій для інших учасників ринку (наразі вже видано понад 30 таких ліцензій).

\section{Модель}

Для розуміння архітектури міжвідомчої взаємодії ми розглядаємо таку базову перелікову модель державних установ:

\begin{enumerate}
    \item Кабінет Міністрів України;
    \item Міністерство освіти і науки України (МОН);
    \item Міністерство охорони здоров'я України (МОЗ);
    \item Міністерство внутрішніх справ України (МВС);
    \item Міністерство закордонних справ України (МЗС);
    \item Міністерство оборони України (МО);
    \item Міністерство юстиції України (Мін'юст);
    \item Міністерство фінансів України (Мінфін);
    \item Міністерство економіки України (Мінекономіки);
    \item Міністерство енергетики України (Міненерго);
    \item Міністерство соціальної політики України (Мінсоцполітики);
    \item Міністерство розвитку громад і територій України (Мінрегіон);
    \item Міністерство цифрової трансформації України (Мінцифра);
    \item Міністерство захисту довкілля та природних ресурсів України (Мінекології).
\end{enumerate}

\newpage
\section{Структурне ядро}

Структурне ядро сучасного цифрового та оптимізованого органу виконавчої влади зазвичай включає три ключові стовпи:
\begin{enumerate}
    \item Управління юстиції внутрішнього контролю;
    \item Департамент інформаційних систем (або виробничо-промислове ІТ-управління);
    \item Управління інституційної трансформації.
\end{enumerate}

Оскільки соціоінформаційні системи повинні залишатися максимально автономними та мати
прозорий керований життєвий цикл, технічне відображення організаційної структури
неодмінно повинно перебувати під загальним контролем профільного міністерства. Найефективніше це реалізовувати у формі окремого
державного підприємства. Такий підхід дає змогу надійно убезпечити ключові інтелектуальні ресурси
й інфраструктуру від неконтрольованого ручного керування, а також обумовлений тим фактом, що ІТ-департаменти міністерств
відповідають за цілісність державних реєстрів та сувору безпеку персональних даних громадян.

Для ефективного керування такою структурою в режимі реального часу ми пропонуємо створити окремий вид
протоколу «ОВВ 2.0». Він стане наступним після розробки класичних НПА (нормативно-правових актів) розширенням до вже існуючого
базового протоколу системи електронного документообігу, що функціонує згідно з постановою №55 
Кабінету Міністрів України. Формальним втіленням цього підходу може стати, наприклад, протокол ДІТ (Державна інституційна трансформація).

У процесі функціонування як операційного документообігу (що генерується структурними
підрозділами міністерств), так і інституційного (який породжується управліннями
або агенціями державної цифрової трансформації) безальтернативно формуються криптографічні зліпки кваліфікованих
електронних підписів (КЕП) відповідальних посадових осіб. Ці сліди дозволяють здійснювати надійний аудит системи — вони перевіряються та досліджуються як внутрішніми
контролюючими органами (відділами аудиту та внутрішніх розслідувань), так і
безпосередньо координуючим наглядовим органом — Міністерством юстиції України.

\newpage
\section{Протокол державної інституційної трансформації}

Цей спеціалізований інформаційний протокол визначає такі базові операції над організаційною структурою відомства та її політиками:

\begin{enumerate}
    \item Створення та ліквідація структурних підрозділів;
    \item Розробка, погодження та модифікація установчих та конституційних документів;
    \item Проєктування та розробка технічних проєктів для структурних підрозділів та їхніх внутрішніх систем;
    \item Збір вимог, вибір та подальше впровадження відповідних інформаційних систем;
    \item Запуск та супровід операційної (безперервної виробничої) діяльності структурних підрозділів.
\end{enumerate}
